% !PRINTED 306
\weberpage{306}

In the proof of this theorem we make the following assumptions, from
some of which we shall later free ourselves again.

1. $f(x)$ has no double or multiple roots.

2. Of the functions $f_\nu(x)$, no two adjacent ones vanish for
the same value of $x$.

3. Likewise, of the functions $F_\nu(x)$, no two adjacent ones vanish
for the same value of $x$.

4. A consequence of 2. is [by formula (3)] that $f_\nu(x)$ and
$F_\nu(x)$ do not vanish for the same value of $x$.

5. For $x = \alpha$ and $x = \beta$ themselves, none of the functions
$f_\nu$, $F_\nu$ is to vanish.

The proof of Theorems II and III is now as follows:

If we let $x$ go, continuously increasing, from $\alpha$ to $\beta$,
then a change in the succession of signs can occur only when one of the
functions $f_\nu$ or $F_\nu$ passes through zero, and it is only a
matter of investigating the effect that such a passage through zero
produces in the various cases. Let us first assume that one of the
middle functions $f$, say $f_\nu$, passes through zero, where
therefore $0 < \nu < n$. According as $f_{\nu+1}$ is positive or negative,
$f_\nu$ will increase or decrease. Thus in the first case
$f_\nu(x - 0)$ will be negative and $f_\nu(x + 0)$ positive, and in the
second case the reverse; moreover it follows from (3) that, when $f_\nu = 0$,
$F_{\nu-1}$ and $F_{\nu+1}$ are positive, and that $F_\nu$ has the
sign opposite to that of $f_{\nu-1} f_{\nu+1}$. We collect
the possibilities that remain after this in tabular form, but omit
from the enumeration those cases which arise from the ones listed by a
simultaneous change of sign of the three functions $f_{\nu-1}$, $f_\nu$,
$f_{\nu+1}$, since these evidently yield no
different results. There remain to us after this only two
cases:

% [mirrored 2026-08-07] The German fragment's lower third was restored after this page's
% translation pass raised the omission (two sign-table displays and the connecting line);
% the displays below are copied character for character from the restored German.
\[
\begin{array}{|c|ccc|}
 & \nu-1, & \nu, & \nu+1 \\ \hline
f(x-0) & + & - & + \\
f(x+0) & + & + & + \\ \hline
F & + & - & +
\end{array}
\qquad
\begin{array}{|ccc|}
\nu-1, & \nu, & \nu+1 \\ \hline
- & - & + \\
- & + & + \\ \hline
+ & + & +
\end{array}.
\]
In these cases we have accordingly
\[
\begin{array}{|c|cc|}
x-0 & VV, & VV \\
x+0 & PV, & PV
\end{array}
\qquad
\begin{array}{|cc|}
PP, & VP \\
VP, & PP
\end{array}.
\]
